\subsection*{两个变化共同决定结果}
面积 $f(x)g(x)$ 随 $x$ 改变时，两条边都可能变化。新增面积包含一条 $g\,\Delta f$ 的带、一条 $f\,\Delta g$ 的带和角上的 $\Delta f\Delta g$。除以输入增量后，前两项分别成为 $gf'$ 与 $fg'$。

角项为什么能消失？可导性意味着两条边的增量都是输入增量的一阶量，所以相乘是二阶量；再除以一次增量，仍会趋零。若没有这项数量级分析，“把角忽略”只是图像猜测。乘积法则计算的是两条独立贡献的和。

\subsection*{串行变化则把倍率相乘}
复合函数的结构不同：输入先经过 $f$，中间输出再经过 $g$。附近小间距先放大约 $f'(x)$ 倍，再放大约 $g'(f(x))$ 倍，所以总倍率相乘。外层导数必须在实际收到的中间值处计算，不能把它误写成 $g'(x)$。

比较 $x^2\sin x$ 和 $\sin(x^2)$。前者同时改变两块因子，因此两条贡献相加；后者把平方结果送入正弦，因此倍率串联。先画计算关系，再选择法则，比看到两个符号便机械判断更可靠。

有时用差商约去中间增量来记忆链式法则很方便，但中间增量也可能为零，严格证明需要用局部线性近似处理。下一章的反向传播会重用同样结构：沿一条计算路径相乘，多条路径汇合时相加。
